% page 212 du fac-simile (image p-211 du scan)
% bloc 160.0 x 242.0 mm ; marge gauche 8.0 mm
\pgc{-12.700mm}{0.000mm}{
\l{}
\l{}
\l{\cel{3}2O4}
\l{}
\l{}
\l{\sou{guelfo}. (\sou{Historio di Italia}). La persono, la partisano qua,}
\l{\cel{3}dum la milito civila lor la Mezepoko, favoris la papo}
\l{\cel{3}(kontraste kun la gibelini, qui favoris la imperiestro}
\l{\cel{3}Germana). - DEFIRS.}
\l{}
\l{\sou{gufo}.\cel{3}(\sou{zool.})\cel{2}Nokt-ucelo de la familio \textquotedbl{}strigi\textquotedbl{}. - L.}
\l{\cel{3}\sou{Otis vulgaris}. - I}
\l{}
\l{\sou{guidar}. - (trans., ad, pri).\cel{2}Montrar la justa voyo ad ulu}
\l{\cel{3}(akompanante lu) qua ne konocas olu, e qua povas hezitar}
\l{\cel{3}od erorar. - EFIS.}
\l{}
\l{\sou{guindo}. (\sou{bot.})\cel{2}Sorto di cerizo, nigra-reda. - L.Cerasus}
\l{\cel{3}\sou{juliana}.}
\l{}
\l{\sou{gul}o.\cel{2}(\sou{mitol. orientala}). Genio qua devoras la korpi di la}
\l{\cel{3}mortinti. - DEFIS.}
\l{}
\l{\sou{+gulfo}. (\sou{geogr.})\cel{2}Maro-parto qua formacas larja konkavajo en}
\l{\cel{3}la tero di litoro. - DEFIRS.}
\l{}
\l{\sou{gulyasho}. Bovo-karno, kun onioni hachita, papriko, alio,}
\l{\cel{3}kumino-grani, majorano, muskado, e c., koquinta dum tri}
\l{\cel{3}o quar hori : a la sauco adjuntesas kelka buliono e farino.}
\l{}
\l{\sou{gumo}. Substanco qua defluas de ula arbori, uzata por glutinar,}
\l{\cel{3}apretar, briligar, o kom dolcigivo en siropi, pasti,}
\l{\cel{3}pastili, e c. - DEFIRS.}
\l{}
\l{\sou{gumiguto}.\cel{2}Gumo ek arboro (L.\cel{1}\sou{garcinia)}\cel{1}di India. - DFIR.}
\l{}
\l{\sou{gumlako}. Gumo ek planti leguminosa di India. - DEFIRS.}
\l{}
\l{\sou{gurdo}. Botelo plata, sur la parto extera di qua aplikesis}
\l{\cel{3}oziero-plektajo, en qua onu varsas vino, brandio, e c.-EF.}
\l{}
\l{\sou{gurmo}. I.\cel{2}Flegmazio di la mukozo nazala, che la yuna kavalo.}
\l{\cel{3}- II. Exantemo di la vizajo e di la pelo haroza (che la}
\l{\cel{3}homo). - F.}
\l{}
\l{\sou{gurmando}. (Persono) qua prizas la dishi bona (bone-saporoza)}
\l{\cel{3}e manjas multo de li. - DEF.}
\l{}
\l{\cel{1}\sou{gurto}. Bendo de ledro o de telo, quan onu tensas por mantenar}
\l{\cel{3}ulo. - DE.}
\l{}
\l{\sou{gustar}. (\sou{trans.}) I. Perceptar la saporo (di ulo). - II. Percep-}
\l{\cel{3}tar la saporo (di ulo) por judikar lo. (\sou{Anke metaf.}).DEFIS.}
\l{}
\l{\sou{guto.}\cel{2}Mikra quanto de liquido qua su separas en la formo di}
\l{\cel{3}globeto. - FS.}
}
